\chapter{Übungen}\label{app:A_Übungen}

\section{Einleitung}

\underline{Aufgabe 1.3.1}

Welche der folgenden physikalischen Grössen sind abgeleitete Grössen? \\
Masse $m$, Geschwindigkeit $v$, Weg $s$, Temperatur $T$, Druck $p$, Kraft $F$, elektrische Spannung $U$,
elektrische Stromstärke $I$, Drehmoment $M$, Zeit $t$, Vorschub $f$, Drehzahl $n$, Fallbeschleunigung $g$ \cite{boge_physik_2008}.

\underline{Aufgabe 1.3.2}

Rechne um in $m$: \\
$4.5 \mu m$; $12.5 \cdot 10^3 nm$; $4.3 \cdot 10^7 pm$; 
$5.8 \cdot 10^4 dm$; $0.04 Gm$; $1.768km$ \cite{boge_physik_2008}

\underline{Aufgabe 1.3.3}

Die Grössengleichung für die Endgeschwindigkeit $v$ eines um die Höhe $h$ frei fallenden Körpers
lautet: $v = \sqrt{2 \cdot g \cdot h}$ . Darin ist $g = 9,81 m s^{-2}$ die Fallbeschleunigung.
Bestimme $v$ mit $h = 20 m$ \cite{boge_physik_2008}. Beweise, dass die Einheit von $v$: $[v] = m s^{-1}$ ist.

\underline{Aufgabe 1.3.4}

Die Grössengleichung für die Leistung $P$ eines Körpers lautet: $P = F \cdot v$.
Bestimme $P$ für eine Kraft $F = 500 N$ und $v = 18 m s^{-1}$ \cite{boge_physik_2008}.

\underline{Aufgabe 1.3.5}

Drücke $P1_{kar}$(2, 2, 1) in Polarkoordinaten $P1_{pol}$ aus.

\underline{Aufgabe 1.3.6}

Drücke $P2_{pol}$(1.414, $\frac{\pi}{4}$, 1) in kartesischen Koordinaten $P2_{kar}$ aus.

\section{Stoffe- und Stoffeigenschaften}

\underline{Aufgabe 2.1.1}

Was ist die Ordnungszahl des Natriumatoms? Wie viele Elektronen hat es?

\underline{Aufgabe 2.1.2}

Warum sind Atome elektrisch neutral?

\underline{Aufgabe 2.1.3}

Wie viele Elektronen hat Silizium (Si) in seiner äusseren Schale?

\underline{Aufgabe 2.1.4}

Mithilfe des Periodensystems, zeichne die Atommodelle (Fokus auf Elektronen, Kern nur symbolisch) für: H, He, Li, Be, B

\underline{Aufgabe 2.2.1}

Zeichne in ein Phasendiagramm, was passiert, wenn bei konstantem Druck ($10 bar$) Wasser von
$-10^\circ C$ auf $400^\circ C$ erwärmt wird. Bestimme und benenne alle durchlaufenen
Übergänge.

\underline{Aufgabe 2.2.2}

Zeichne in ein Phasendiagramm, was passiert, wenn der Druck gleichmässig von $0 bar$ auf $221 bar$ erhöht
und zugleich das Wasser von $374^\circ C$ auf $-10^\circ C$ abgekühlt wird. Bestimme und benenne alle durchlaufenen
Übergänge.

\section{Statik}

\underline{Aufgabe 3.1.1}

\begin{minipage}{.58\textwidth}
    Zerlege die Kraft $F$ in seine Komponenten $F_x$ und $F_y$. \newline
    $F = 1kN$ \newline
    $\displaystyle \alpha = \frac{\pi}{4}$
\end{minipage}
\hspace{.1\textwidth}
\begin{minipage}{.3\textwidth}
    \begin{figure}[H]    
        \centering          %     left botm right top
        \includegraphics[clip, trim=0cm 0cm 0cm 0cm,width=\textwidth]{kraftImKartesischen.png} %.6 stands for 0.6 x textwidth
        \label{fig:311}
    \end{figure}
\end{minipage}

\underline{Aufgabe 3.2.1}

\begin{minipage}{.58\textwidth}
    Mit dem Schlüssel wird eine Schraube gelöst. Wie gross ist das
    resultierende Moment bezüglich der Achse der Schraube (Punkt 0) \newline
    $F_1 = 100N$ \newline
    $F_2 = 80N$
\end{minipage}
\hspace{.1\textwidth}
\begin{minipage}{.3\textwidth}
    \begin{figure}[H]    
        \centering          %     left botm right top
        \includegraphics[clip, trim=0cm 0cm 0cm 0cm,width=\textwidth]{aufgabe_3_2_1.png} %.6 stands for 0.6 x textwidth
        \label{fig:321}
    \end{figure}
\end{minipage}

\underline{Aufgabe 3.2.2}

\begin{minipage}{.58\textwidth}
    Drei Kräfte wirken auf den Körper. \newline
    Das Mass $a$ ist $2m$. \newline
    Zu bestimmen ist das Moment der Kräfte bezüglich dem Punkt A.
\end{minipage}
\hspace{.1\textwidth}
\begin{minipage}{.3\textwidth}
    \begin{figure}[H]    
        \centering          %     left botm right top
        \includegraphics[clip, trim=0cm 0cm 0cm 0cm,width=\textwidth]{aufgabe_3_2_2.png} %.6 stands for 0.6 x textwidth
        \label{fig:322}
    \end{figure}
\end{minipage}

\underline{Aufgabe 3.2.1.1}

\begin{minipage}{.58\textwidth}
    Die auf die Wellen des Getriebes wirkenden Momente sind 
    $M_1 = 500Nm, M_2 = 60Nm, M_3 = 40Nm$.\newline
    Gesucht: Die Kräfte in A und B, so dass das Moment kompensiert wird 
    ($l = 350mm$).
\end{minipage}
\hspace{.1\textwidth}
\begin{minipage}{.3\textwidth}
    \begin{figure}[H]    
        \centering          %     left botm right top
        \includegraphics[clip, trim=0cm 0cm 0cm 0cm,width=\textwidth]{aufgabe_3_2_1_1.png} %.6 stands for 0.6 x textwidth
        \label{fig:3211}
    \end{figure}
\end{minipage}

\underline{Aufgabe 3.2.2.1}

\begin{minipage}{.58\textwidth}
    Die Lage des Schwerpunktes der zusammengesetzten Fläche ist zu bestimmen. 
    Das vorgegebene Koordinatensystem ist zu verwenden.
\end{minipage}
\hspace{.1\textwidth}
\begin{minipage}{.3\textwidth}
    \begin{figure}[H]    
        \centering          %     left botm right top
        \includegraphics[clip, trim=0cm 0cm 0cm 0cm,width=\textwidth]{aufgabe_3_2_2_1.png} %.6 stands for 0.6 x textwidth
        \label{fig:3221}
    \end{figure}
\end{minipage}

\underline{Aufgabe 3.3.1.1}

\begin{minipage}{.58\textwidth}
    2 Kräfte greifen nach Skizze an einem Haltestab an. Mit welcher resultierenden 
    Kraft wird am Stab gezogen und mit welcher gebogen? \newline
    $F_{1} = 10kN$ \newline
    $F_2 = 6kN$
\end{minipage}
\hspace{.1\textwidth}
\begin{minipage}{.3\textwidth}
    \begin{figure}[H]    
        \centering          %     left botm right top
        \includegraphics[clip, trim=0cm 0cm 0cm 0cm,width=\textwidth]{aufgabe_3_3_1_1.png} %.6 stands for 0.6 x textwidth
        \label{fig:3311}
    \end{figure}
\end{minipage}

\underline{Aufgabe 3.3.1.2}

\begin{minipage}{.58\textwidth}
    Alle Seilkräfte an der Halterung und die Winkel $\beta$ und $\gamma$ sind vorgegeben. 
    In welcher Richtung muss die Kraft $F_1$ wirken, so dass die Halterung nur 
    vertikal belastet wird? \newline
    $F_1 = 3kN$ \newline
    $F_2 = 2.5kN, \gamma = 80^{\circ}$ \newline
    $F_3 = 1.8kN, \beta = 30^{\circ}$
\end{minipage}
\hspace{.1\textwidth}
\begin{minipage}{.3\textwidth}
    \begin{figure}[H]    
        \centering          %     left botm right top
        \includegraphics[clip, trim=0cm 0cm 0cm 0cm,width=\textwidth]{aufgabe_3_3_1_2.png} %.6 stands for 0.6 x textwidth
        \label{fig:3312}
    \end{figure}
\end{minipage}

\underline{Aufgabe 3.3.1.3}

\begin{minipage}{.58\textwidth}
    Zwei Schlepper ziehen geradlinig wie abgebildet ein Schiff. Die Seilkraft 
    $F2 = 50kN$ und die beiden Winkel werden gemessen. \newline
    Zu bestimmen sind die Seilkraft $F_1$ und die am Schiff angreifende 
    Resultierende.
\end{minipage}
\hspace{.1\textwidth}
\begin{minipage}{.3\textwidth}
    \begin{figure}[H]    
        \centering          %     left botm right top
        \includegraphics[clip, trim=0cm 0cm 0cm 0cm,width=\textwidth]{aufgabe_3_3_1_3.png} %.6 stands for 0.6 x textwidth
        \label{fig:3313}
    \end{figure}
\end{minipage}

\underline{Aufgabe 3.3.2.1}

\begin{minipage}{.58\textwidth}
    Die Resultierende des Kräftesystems ist nach Betrag und Lage zu bestimmen. \newline
    Das Ergebnis ist mit einem von A verschiedenen Pol zu kontrollieren.
\end{minipage}
\hspace{.1\textwidth}
\begin{minipage}{.3\textwidth}
    \begin{figure}[H]    
        \centering          %     left botm right top
        \includegraphics[clip, trim=0cm 0cm 0cm 0cm,width=\textwidth]{aufgabe_3_3_2_1.png} %.6 stands for 0.6 x textwidth
        \label{fig:3321}
    \end{figure}
\end{minipage}

\underline{Aufgabe 3.3.2.2}

\begin{minipage}{.58\textwidth}
    Der Träger ist in A und B gelagert. \newline
    Die in B angreifende Kraft $B$ ist so zu bestimmen, dass das resultierende
    Moment bezüglich A verschwindet \newline
    Welche Bedeutung hat diese Bedingung?
\end{minipage}
\hspace{.1\textwidth}
\begin{minipage}{.3\textwidth}
    \begin{figure}[H]    
        \centering          %     left botm right top
        \includegraphics[clip, trim=0cm 0cm 0cm 0cm,width=\textwidth]{aufgabe_3_3_2_2.png} %.6 stands for 0.6 x textwidth
        \label{fig:3322}
    \end{figure}
\end{minipage}

\underline{Aufgabe 3.3.3.1}

\begin{minipage}{.58\textwidth}
    An einem Hebel greifen 3 Kräfte an. \newline
    Es sind die auf die Welle 0 wirkende resultierende Kraft nach Betrag
    und Richtung und das resultierende Moment zu berechnen.
\end{minipage}
\hspace{.1\textwidth}
\begin{minipage}{.3\textwidth}
    \begin{figure}[H]    
        \centering          %     left botm right top
        \includegraphics[clip, trim=0cm 0cm 0cm 0cm,width=\textwidth]{aufgabe_3_3_3_1.png} %.6 stands for 0.6 x textwidth
        \label{fig:3331}
    \end{figure}
\end{minipage}

\underline{Aufgabe 3.3.3.2}

\begin{minipage}{.58\textwidth}
    An einem dreieckförmigen Körper greifen gem. Skizze drei Kräfte und ein Moment an. \newline
    Zu bestimmen sind die resultierende Kraft nach Betrag und Richtung 
    sowie das resultierende Moment bez. des Schwerpunktes S.
\end{minipage}
\hspace{.1\textwidth}
\begin{minipage}{.3\textwidth}
    \begin{figure}[H]    
        \centering          %     left botm right top
        \includegraphics[clip, trim=0cm 0cm 0cm 0cm,width=\textwidth]{aufgabe_3_3_3_2.png} %.6 stands for 0.6 x textwidth
        \label{fig:3332}
    \end{figure}
\end{minipage}

\underline{Aufgabe 3.3.3.3}

\begin{minipage}{.58\textwidth}
    An der angeschweissten Lasche greift die Kraft $F$ an.
    Zur Berechnung der Festigkeit des Trägers ist die 
    Last auf eine gleichwertige in der Trägerachse (Punkt $M$) zu ersetzen.
\end{minipage}
\hspace{.1\textwidth}
\begin{minipage}{.3\textwidth}
    \begin{figure}[H]    
        \centering          %     left botm right top
        \includegraphics[clip, trim=0cm 0cm 0cm 0cm,width=\textwidth]{aufgabe_3_3_3_3.png} %.6 stands for 0.6 x textwidth
        \label{fig:3333}
    \end{figure}
\end{minipage}

\underline{Aufgabe 3.3.3.4}

\begin{minipage}{.6\textwidth}
    Am Hebel greifen die horizontale Kraft $F_A$
    und die Kraft $F_B$ in Richtung $\alpha$ an.
    
    Die beiden Kräfte sollen so bestimmt werden,
    dass am Vierkant die resultierende Kraft
    $F = 200N$ und das resultierende Moment
    $M = 80Nm$ wirken. (8)

\end{minipage}
\hspace{.04\textwidth}
\begin{minipage}{.35\textwidth}
    \begin{figure}[H]   
        \includegraphics[width=\textwidth]{aufgabe_3_3_3_4.png} %.6 stands for 0.6 x textwidth
    \end{figure}
\end{minipage}

Tipp: Beginne mit der Summe der Momente und Summe der Kräfte in y um $F_A$ zu berechnen.

\underline{Aufgabe 3.4.1}

\begin{minipage}{.5\textwidth}
    \begin{figure}[H]    
        \centering          %     left botm right top
        \includegraphics[clip, trim=0cm 0cm 0cm 0cm,width=\textwidth]{aufgabe_3_4_1.png} %.6 stands for 0.6 x textwidth
    \end{figure}
\end{minipage}
\begin{minipage}{.5\textwidth}
    \begin{figure}[H]    
        \centering          %     left botm right top
        \includegraphics[clip, trim=0cm 0cm 0cm 0cm,width=\textwidth]{aufgabe_3_4_1_right.png} %.6 stands for 0.6 x textwidth
    \end{figure}
\end{minipage}

\section{Kinematik}

\underline{Aufgabe 4.2.1}

Ein Förderband mit einer Neigung von $50^\circ$ zur Waagerechten überwindet einen Höhenunterschied
von $30m$ in $4min$.

Bestimme die Geschwindigkeit des Fördergutes auf der schiefen Ebene \cite{boge_physik_2008}.

\underline{Aufgabe 4.2.2}

\begin{minipage}{.58\textwidth}
    Die Trommel A wickelt das Seil mit konstanter
    Geschwindigkeit $v_A = 4 ms^{-1}$ auf, die Trommel B
    wickelt das Seil mit konstanter Geschwindigkeit
    $v_B = 3 ms^{-1}$ ab. Zu bestimmen sind:

    a) die Bewegungsgleichung für die Last D

    b) Lage und Geschwindigkeit der Last D
    für $t = 5s$, wenn für $t = 0s$ $s = 0m$ ist

\end{minipage}
\hspace{.1\textwidth}
\begin{minipage}{.3\textwidth}
    \begin{figure}[H]    
        \centering          %     left botm right top
        \includegraphics[clip, trim=0cm 0cm 0cm 0cm,width=\textwidth]{aufgabe_4_2_2.png} %.6 stands for 0.6 x textwidth
        \label{fig:421}
    \end{figure}
\end{minipage}

\underline{Aufgabe 4.2.3}

Die Geschwindigkeit eines Körpers nimmt in $5s$ um $40kmh^{-1}$ ab.
Wie gross sind die Geschwindigkeitsänderung $\Delta v$ in $ms^{-1}$ und die Verzögerung $a$ in $ms^{-2}$ \cite{boge_physik_2008}?

\underline{Aufgabe 4.2.4}

Skizziere das v-t-Diagramm

a) für den freien Fall
b) für den senkrechten Wurf (nach oben) bis zur Rückkehr des Körpers zur Abwurfstelle \cite{boge_physik_2008}.

\underline{Aufgabe 4.2.5}

Bestimme die Fallhöhe $h$, die ein Körper haben müsste, um beim freien Fall die 
Endgeschwindigkeit $v_e = 30 ms^{-1}$ zu bekommen \cite{boge_physik_2008}.

\underline{Aufgabe 4.2.6}

\begin{minipage}{.58\textwidth}
    Ein Motorrad startet aus dem Stillstand und beschleunigt konstant, 
    bis es nach 100 m eine Geschwindigkeit von $108kmh^{-1}$ erreicht.

    a) Wie gross ist die Beschleunigung?

    b) Wie lange braucht es dafür?

    c) Bei einer anschliessenden konstanten Verzögerung von $a=-7ms^{-2}$,
    nach welcher Strecke kommt das Motorrad wieder zum stehen?

\end{minipage}
\hspace{.1\textwidth}
\begin{minipage}{.3\textwidth}
    \begin{figure}[H]    
        \centering          %     left botm right top
        \includegraphics[clip, trim=0cm 0cm 0cm 0cm,width=\textwidth]{aufgabe_4_2_1.jpg}
        \label{fig:422}
    \end{figure}
\end{minipage}

\underline{Aufgabe 4.3.1}

\begin{minipage}{.58\textwidth}
    Ein Auto fährt mit einer Geschwindigkeit von $60kmh^{-1}$ von einer Klippe.

    a) Welche Geschwindigkeit hat das Fahrzeug beim Aufprall auf das $30m$ tiefer gelegte
    Flussbett?

    b) Wie weit von der Klippe entfernt schlägt das Fahrzeug auf ($x$-Richtung)? 

\end{minipage}
\hspace{.1\textwidth}
\begin{minipage}{.3\textwidth}
    \begin{figure}[H]    
        \centering          %     left botm right top
        \includegraphics[clip, trim=0cm 0cm 0cm 0cm,width=\textwidth]{aufgabe_4_3_1.png}
        \label{fig:431}
    \end{figure}
\end{minipage}

\underline{Aufgabe 4.4.1}

Bestimme die (durchschnittliche) Winkelgeschwindigkeit

a) für den kleinen Zeiger einer Uhr

b) für den großen Zeiger.

\underline{Aufgabe 4.4.2}

Eine mit $n_1 = 5000min^{-1}$ laufende Motorwelle wird innerhalb 
von $10 s$ auf $n_2 = 3000 min^{-1}$ gleichmässig abgebremst. Bestimme: 

a) die Winkelverzögerung $\alpha$

b) die Anzahl $z$ der Umdrehungen während des Bremsvorgangs

\underline{Aufgabe 4.4.3}

Eine Schleifscheibe von $d = 400mm$ Durchmesser wird aus dem Stillstand heraus in $20s$ gleich-
mäßig auf $v_u = 30 ms^{-1}$ beschleunigt \cite{boge_physik_2008}.
Bestimme: 

a) die Enddrehzahl $n_t$ , 

b) die Endwinkelgeschwindigkeit $\omega_t$

c) die Anzahl $z$ der Umdrehungen beim Beschleunigen

d) die Winkelbeschleunigung $\alpha$

e) die Tangentialbeschleunigung $a_T$ eines Korns auf dem Scheibenumfang.

\underline{Aufgabe 4.4.4}

\begin{minipage}{.58\textwidth}
    Ein Satellit befindet sich auf einer Erdumlaufbahn mit einer Geschwindigkeit $v$ von
    $7.8 kms^{-1}$. Die Distanz zu der Erde ist konstant. Es wird angenommen, dass die 
    Gravitation noch ca. $g = 9.2ms^{-2}$ beträgt. Die Erde hat einen
    Durchmesser von ca. $d = 12700km$.

    Wie gross muss die Distanz $s$ sein?

\end{minipage}
\hspace{.1\textwidth}
\begin{minipage}{.3\textwidth}
    \begin{figure}[H]    
        \centering          %     left botm right top
        \includegraphics[clip, trim=0cm 0cm 0cm 0cm,width=\textwidth]{aufgabe_4_4_1.png}
        \label{fig:441}
    \end{figure}
\end{minipage}

\section{Dynamik}

\underline{Aufgabe 5.1.1}

\begin{minipage}{.58\textwidth}
Dackel Waldemar mit seinem Kampfgewicht von $4kg$ steht nichtsahnend in einem mit konstanter Geschwindigkeit geradeaus
fahrenden $100t$ schweren Bahnwagen, als dieser plötzlich über eine Weiche fährt und mit einer Beschleunigung von $1ms^{-2}$
quer zur Fahrtrichtung das Gleis wechselt. Zeichne die folgenden Kraftpfeile mit Betrag und Richtung in die Skizzen ein.

a) $\overrightarrow{F}_{W \rightarrow S}$, Kraft Wagen auf Schiene

b) $\overrightarrow{F}_{S \rightarrow W}$, Kraft Schiene auf Wagen

c) $\overrightarrow{F}_{W \rightarrow D}$, Kraft Wagen auf Dackel

d) $\overrightarrow{F}_{D \rightarrow W}$, Kraft Dackel auf Wagen

e) Welche Beschleunigung erfährt der Dackel in Bezug auf das
System des Wagens?

f) Welche Beschleunigung erfährt der Dackel in Bezug auf das
System der Schiene?

\end{minipage}
\hspace{.1\textwidth}
\begin{minipage}{.3\textwidth}
    \begin{figure}[H]    
        \centering          %     left botm right top
        \includegraphics[clip, trim=0cm 0cm 0cm 0cm,width=\textwidth]{aufgabe_dackel.png}
        \label{fig:dackel}
    \end{figure}
\end{minipage}

\underline{Aufgabe 5.1.2}

\begin{minipage}{.58\textwidth}
    Auf dem einen Ende eines 1 m langen Brettes liegt ein Holzklotz mit
    Haftreibungszahl $\mu_{HR} = 0.8$ und Gleitreibungszahl $\mu_{GR} = 0.6$. 
    
    a) Wie hoch kann man das Brett auf der Seite anheben, bis der Klotz ins
    Rutschen gerät? 
    
    b) Wie schnell ist er dann am unteren Ende?

\end{minipage}
\hspace{.1\textwidth}
\begin{minipage}{.3\textwidth}
    \begin{figure}[H]    
        \centering          %     left botm right top
        \includegraphics[clip, trim=0cm 0cm 0cm 0cm,width=\textwidth]{aufgabe_5_1_1.png}
        \label{fig:511}
    \end{figure}
\end{minipage}

\underline{Aufgabe 5.1.3}

\begin{minipage}{.58\textwidth}
    Berechne jeweils die Beschleunigung der beiden
    reibungsfrei gelagerten und über eine Schnur auf einer
    ebenfalls reibungsfreien Rolle verbundenen Körper,
    wenn sie sich unter dem Einfluss der Gravitationskraft
    anfangen zu bewegen.

\end{minipage}
\hspace{.1\textwidth}
\begin{minipage}{.3\textwidth}
    \begin{figure}[H]    
        \centering          %     left botm right top
        \includegraphics[clip, trim=0cm 0cm 0cm 0cm,width=\textwidth]{aufgabe_5_1_2.png}
        \label{fig:512}
    \end{figure}
\end{minipage}

\underline{Aufgabe 5.2.1}

\begin{minipage}{.58\textwidth}
    Ein Kletterer ($80kg$) stürzt 5m über der letzten Zwischensicherung und
    fällt frei ins elastische Seil ($A= 100 mm^2$, $E=300Nmm^{-2}$). Das Seil ist
    vor dem Sturz $30m$ lang (straff) und der sichernde Partner hält das
    Seil fix. Reibungsverluste sollen vernachlässigt werden. Gesucht sind:
    
    a) Maximale Geschwindigkeit des Kletterers
    
    b) Fallhöhe h
    
    c) Maximale Seilkraft
    
    d) Maximale Verzögerung des Kletterers

\end{minipage}
\hspace{.1\textwidth}
\begin{minipage}{.3\textwidth}
    \begin{figure}[H]    
        \centering          %     left botm right top
        \includegraphics[clip, trim=0cm 0cm 0cm 0cm,width=\textwidth]{aufgabe_5_2_1.png}
        \label{fig:521}
    \end{figure}
\end{minipage}

\underline{Aufgabe 5.2.2}

Ein Auto ($m=950kg$) wird in $4s$ von $v_1 = 50 kmh^{-1}$
auf $v_2 = 90 kmh^{-1}$ beschleunigt.

a) Welchen Weg legt es dabei zurück.

b) Wie groß ist die Beschleunigungsarbeit?

c) Welche Geschwindigkeit hätte der Wagen mit der gleichen Beschleunigungsarbeit
erreicht, wenn er aus dem Stand beschleunigt wäre?

\underline{Aufgabe 5.2.3}

Die Kugel eines Fadenpendels ($l=1.2m$) wird aus der Gleichgewichtslage um $60 cm$
ausgelenkt. Welche Geschwindigkeit muss die Kugel zusätzlich erhalten, damit sie nach dem
Zurückschwingen eine Kreisbahn beschreibt?

\underline{Aufgabe 5.3.1}

Urs wiegt $40 kg$ und rollt mit $2 ms^{-1}$ auf seinem $2 kg$ schweren Skateboard auf die Mädchen zu. 
Jetzt springt er nach hinten ab, so dass er auf $v = 1 ms^{-1}$ abbremst und lässig zum Stehen kommt. 
Wie schnell schießt das Skateboard jetzt in Richtung auf die kleine Lisa davon?

\underline{Aufgabe 5.3.2}

Der $60 kg$ schwere Max stösst sich mit $200 N$ während $0.5 s$ von dem $200 kg$ schweren Tretboot ab, um an das $1 m$
entfernte Ufer zu springen. Wie schnell wird er dabei? Wie weit hat sich das Boot in der halben Sekunde schon vom Ufer
fortbewegt?

\underline{Aufgabe 5.3.3}

\begin{minipage}{.58\textwidth}
    Ein Kerbschlagversuch wird mit einem instrumentierten Pendelschlagwerk durchgeführt, so dass
    das nebenstehende Kraft-Zeit Diagramm aufgenommen werden kann. Der vertikale Abstand
    Probe - Pendelmasse beträgt $H = 1.5 m$, die Pendelmasse $m = 11 kg$.
    
    Wie gross ist die Schlagarbeit?

\end{minipage}
\hspace{.1\textwidth}
\begin{minipage}{.3\textwidth}
    \begin{figure}[H]    
        \centering          %     left botm right top
        \includegraphics[clip, trim=0cm 0cm 0cm 0cm,width=\textwidth]{aufgabe_pendelschlagwerk.png}
        \label{fig:532}
    \end{figure}
\end{minipage}

\underline{Aufgabe 5.3.4}

\begin{minipage}{.58\textwidth}
    Die 2 Massen $J_A$ und $J_B$ mit den Drehzahlen $n_A$
    und $n_B$ rotieren frei. Wie gross ist die Drehzahl,
    wenn sie gekoppelt werden?

\end{minipage}
\hspace{.1\textwidth}
\begin{minipage}{.3\textwidth}
    \begin{figure}[H]    
        \centering          %     left botm right top
        \includegraphics[clip, trim=0cm 0cm 0cm 0cm,width=\textwidth]{aufgabe_5_2_2.png}
        \label{fig:522}
    \end{figure}
\end{minipage}

\underline{Aufgabe 5.3.5}

Ein $50 g$ schwerer Wagen prallt mit $1 ms^{-1}$ auf einen $400 g$ schweren ruhenden Wagen. 
Wie schnell sind die Wagen nach dem Stoss?

a) im elastischen Fall?

b) Im teilelastische Fall, wenn 30\% der Bewegungsenergie in Wärme umgewandelt werden?

c) im unelastischen Fall? Berechne in diesem Fall den prozentualen Energieverlust.

\underline{Aufgabe 5.3.6}

\begin{minipage}{.58\textwidth}
    Eine Billardkugel trifft mit $1ms^{-1}$ reibungsfrei eine zweite gleich schwere Kugel, die
    um den halben Kugeldurchmesser versetzt zur Bewegungsrichtung ruht. (siehe
    Skizze). Die beiden Kugeln bewegen sich im Fall eines vollkommen elastischen
    Stosses in einem rechten Winkel voneinander fort.

    Bestimme die Richtungen und Beträge der Geschwindigkeiten der Kugeln nach
    dem Stoss für diesen Fall.
\end{minipage}
\hspace{.1\textwidth}
\begin{minipage}{.3\textwidth}
    \begin{figure}[H]    
        \centering          %     left botm right top
        \includegraphics[clip, trim=0cm 0cm 0cm 0cm,width=\textwidth]{aufgabe_5_3_6.png}
        \label{fig:536}
    \end{figure}
\end{minipage}

\section{Schwingungen und Wellen}

\underline{Aufgabe 6.1.1}

Bei einer Schwingung werden in $20 s$ 10 Schwingungen gezählt. 
Der Betrag der Amplitude von $5.0 cm$ bleibt konstant. 
Berechnen Sie Frequenz, Periodendauer
(Schwingungsdauer), Kreisfrequenz sowie maximale Geschwindigkeit und Beschleunigung.

\underline{Aufgabe 6.1.2}

\begin{minipage}{.58\textwidth}
    In der folgenden Skizze finden Sie das Weg-Zeit-Diagramm für eine 
    ungedämpfte harmonische Bewegung.

    a) Bestimme aus diesem Diagramm und den angegebenen Daten

    a1) die Amplitude $\hat{y}$

    a2) die Schwingungsdauer $T_0$

    a3) die Eigenfrequenz $f_0$

    a4) die Eigenkreisfrequenz $\omega_0$

    a5) den Nullphasenwinkel $\phi_0$

    b) Stelle die Bewegungsgleichung $y(t)$ für die gezeichnete ungedämpfte
    harmonische Schwingung auf

    c) Berechne die Geschwindigkeit v0 = v(0) für den Zeitpunkt t = 0.
\end{minipage}
\hspace{.1\textwidth}
\begin{minipage}{.3\textwidth}
    \begin{figure}[H]    
        \centering          %     left botm right top
        \includegraphics[clip, trim=0cm 0cm 0cm 0cm,width=\textwidth]{aufgabe_6_1_2.png}
        \label{fig:612}
    \end{figure}
\end{minipage}

\underline{Aufgabe 6.1.3}

\begin{minipage}{.58\textwidth}
    Ein Fadenpendel schwingt mit einer Auslenkungsamplitude von $A = \frac{\pi}{4}$.
    Die Eigenkreisfrequenz ist $\omega_0 = \pi rads^{-1}$. Die Masse der 
    Kugel beträgt $m = 0.2kg$.

    a) Wie lange ist der Faden $l$?

    b) Stelle die Gleichung auf für $\varphi(t)$. (Ohne Phasenverschiebung)

    c) $0.1s$ nach dem Nulldurchgang, wie ist der Winkel $\varphi(0.1)$?

    d) Wie gross sind Winkelgeschwindigkeit und Winkelbeschleunigung nach $0.1s$?

\end{minipage}
\hspace{.1\textwidth}
\begin{minipage}{.3\textwidth}
    \begin{figure}[H]    
        \centering          %     left botm right top
        \includegraphics[clip, trim=0cm 0cm 0cm 0cm,width=\textwidth]{fadenpendel_aufgabe.png}
        \label{fig:61zwischendrin}
    \end{figure}
\end{minipage}

\underline{Aufgabe 6.1.4}

\begin{minipage}{.58\textwidth}
    Eine kleine Kugel der Masse $m = 100 g$ fällt aus der Höhe $h = 20 cm$ auf eine
    entspannte und masselose Feder mit der Federkonstanten $k = 20 N/m$
    Die Kugel haftet auf der Feder und führt harmonische Schwingungen aus. Von
    Reibungseinflüssen ist abzusehen.
\end{minipage}
\hspace{.1\textwidth}
\begin{minipage}{.3\textwidth}
    \begin{figure}[H]    
        \centering          %     left botm right top
        \includegraphics[clip, trim=0cm 0cm 0cm 0cm,width=\textwidth]{aufgabe_6_1_3.png}
        \label{fig:613}
    \end{figure}
\end{minipage}

a) Um welche maximale Strecke $y = y_{max}$ wird 
die entspannte Schraubenfeder zusammengedrückt?

b) Mit welcher Frequenz $f_0$ schwingt das Feder-Massesystem?

c) Um welche Gleichgewichtslage $y = y_0$ erfolgt die Schwingung?

d) Wie gross ist die Amplitude der Schwingung $\hat{y}$

e) Wie lauten die Anfangsbedingungen $y(0)$ und $v(0)$ im Zeitpunkt 
$t = 0$ des Auftreffens der Kugel auf die entspannte Feder?

f) Skizzieren Sie den qualitativen Verlauf der Schwingung $y(t)$ im y-t Diagramm.

g) Berechnen Sie den Nullphasenwinkel der harmonischen Schwingung y(t) aus den
Anfangsbedingungen.

h) Geben Sie die Formel für die Lösung y(t) an und legen Sie dabei die in der
Abbildung festgelegte y-Achse zu Grunde.

\underline{Aufgabe 6.1.5}

\begin{minipage}{.78\textwidth}
    Ein Feder-Masse-System ($m = 0.95 kg$, $k = 90 Nm^{-1}$) hat bei einer
    geschwindigkeitsproportionalen Dämpfung eine Schwingungsdauer von $T_d = 0.66 s$.

    a) Welche Eigenkreisfrequenz
    $\omega_0$ hat das ungedämpfte System?
    Welchen Dämpfungsgrad $D$ hat das gedämpfte System?

    b) Wie gross ist die Auslenkung $y(t)$ des Körpers nach zwei
    Schwingungsperioden, wenn die Anfangsauslenkung $8 cm$ war?

\end{minipage}
\hspace{.1\textwidth}
\begin{minipage}{.1\textwidth}
    \begin{figure}[H]    
        \centering          %     left botm right top
        \includegraphics[clip, trim=0cm 0cm 0cm 0cm,width=\textwidth]{aufgabe_6_1_5.png} %.6 stands for 0.6 x textwidth
        \label{fig:615}
    \end{figure}
\end{minipage}

\underline{Aufgabe 6.2.1}

Welche Wellenlänge $\lambda$ hat eine Rundfunkwelle von $f = 50 MHz$ bei 
einer Ausbreitungsgeschwindigkeit $c = 300 000 kms^{-1}$?

\underline{Aufgabe 6.2.2}

Die Ausbreitungsgeschwindigkeit einer Querwelle mit der Amplitude $A = 20 cm$ beträgt
$c = 100 ms^{-1}$, die Frequenz $f = 0,1 Hz$.
Bestimme:

a) die Wellenlänge $ \lambda$

b) den Zeitabschnitt $\Delta t$, nach dem ein Teilchen im 
Abstand $\Delta x = 400 m$ vom Ursprung zu schwingen beginnt

c) die Auslenkung dieses Teilchen nach $t_0 = 1 min$.

\underline{Aufgabe 6.2.3}

\begin{minipage}{.58\textwidth}
    Eine laufende Querwelle hat die Amplitude $A = 20 mm$ und
    die Wellenlänge $\lambda = 50 cm$.
    \vspace{8pt}
    Bestimme die Auslenkung $y$ eines Teilchens $P$, das
    $\Delta x = 30 cm$ vom Anfangsteilchen entfernt ist, wenn das
    Anfangsteilchen $A$ gerade eine volle Periode ausgeführt hat.
\end{minipage}
\hspace{.1\textwidth}
\begin{minipage}{.3\textwidth}
    \begin{figure}[H]    
        \centering          %     left botm right top
        \includegraphics[clip, trim=0cm 0cm 0cm 0cm,width=\textwidth]{aufgabe_6_2_3.png} %.6 stands for 0.6 x textwidth
        \label{fig:623}
    \end{figure}
\end{minipage}

\underline{Aufgabe 6.2.4}

Bestimme die momentane Geschwindigkeit $v_y$ und die momentane Beschleunigung $a_y$ des Teilchens
aus der Aufgabe 6.2.3, wenn die Ausbreitungsgeschwindigkeit $c = 5 ms^{-1}$ beträgt.